\section{Conclusions}
\label{sec:conclusions}

HLG provides a small, fully-verifiable agentic benchmark that exercises long-horizon
planning, irreversible side-effect modeling, and dead-end avoidance on a discrete
spatial domain that frontier LLMs have not been heavily trained on. The benchmark is
deliberately compact (34 levels, ~150 lines of public Python SDK, one engine module per
concern), to make it tractable for independent re-implementation and to keep the
evaluation cost low enough to run as a regression test on every new model release.

Three open problems stand out:

\begin{itemize}[leftmargin=1.5em]
    \item \textbf{Dead-state reasoning.} Models do not yet appear to distinguish
      \texttt{[LOSS]} from \texttt{[DEAD]} in their reasoning chains. Eliciting that
      distinction may require explicit backward-reachability heuristics in the system
      prompt, or training signals that reward identifying the earlier branch point in a
      dead-attempt trajectory.
    \item \textbf{Weak-tile orientation reasoning.} The interaction between block
      orientation and tile fragility is conceptually simple but operationally subtle;
      it is a microcosm of physical-affordance reasoning, and current models struggle
      with it.
    \item \textbf{Split-mode merge planning.} Split mode requires planning over two
      independently controlled bodies that must be reunited; the merge condition
      depends on adjacency at the time of merge. This is the closest HLG approximation
      to multi-agent coordination over shared state.
\end{itemize}

We anticipate v0.2 of HLG will add controlled human trials, a public REST surface
implementing \texttt{hlg-openapi.yaml}, provider adapters under \texttt{hlg.providers},
and an expanded private level set.
